
%%%%%%%%%%%%%%%%%%%%%%%%%%%%%%%%%%%%
% Recap/Agenda 
%%%%%%%%%%%%%%%%%%%%%%%%%%%%%%%%%%%%
% TODO better formatting
\begin{frame}
    \frametitle{Module 5: Linear regression}
    \begin{itemize}
        \item \hl{Previously: }Model Selection (Chapter 9.2)
        \item \hl{This time: }Checking model conditions using graphs (Chapter 9.3)
        \item \hl{Reading: }
        \item \hl{Deadlines/Announcements: }
    \end{itemize}
    
\end{frame}

%%%%%%%%%%%%%%%%%%%%%%%%%%%%%%%%%%%%
% Learning objectives:
%%%%%%%%%%%%%%%%%%%%%%%%%%%%%%%%%%%%
\begin{frame}
    \frametitle{Learning Objectives}
    \begin{itemize}
        \item \textbf{M5, LO3: Fit and Interpret Linear Models Using Least Squares:} Fit the intercept and slope of a linear model to data using the least squares method, interpret the fitted values, and use the model to predict responses to new inputs.
        \item \textbf{M5, LO4: Explain and Assess Model Fit:} Explain the least squares fitting procedure, assess whether the fit is unduly influenced by any particular points, and distinguish between interpolation and extrapolation in predictions.
        \item \textbf{M5, LO5: Perform Inference for Regression Coefficients:} Use fit summary and parameter estimates (e.g., $\hat{\beta}_1$ and $\hat{\sigma}^2$) to perform hypothesis tests or construct confidence intervals for the slope, and interpret the results.
    \end{itemize}
\end{frame}
    
